\section{Metodologia}

Este trabalho propõe um método de inspeção visual ativa para o diagnóstico de doenças
em tomateiro utilizando Aprendizado por Reforço Profundo. O diagnóstico é modelado como
uma sequência de decisões: um agente controla uma janela de atenção, desloca e
redimensiona essa janela sobre a imagem e, ao final da inspeção, emite uma decisão
entre \textit{GrayMold}, \textit{Viral} e \textit{Wilt}. O ambiente foi implementado em
Gymnasium, e o agente foi treinado com Proximal Policy Optimization (PPO)
\cite{schulman2017ppo}.

\subsection{Base de Dados}

Foi utilizado o \textit{Tomato Disease Dataset}, composto por 1.026 imagens e 3.167
caixas delimitadoras anotadas em formato Pascal VOC. As anotações indicam regiões
sintomáticas em folhas, frutos e caules. O conjunto foi dividido de forma estratificada
em treino, validação e teste.

\subsection{Formulação como MDP}

O problema foi modelado como um Processo de Decisão de Markov (MDP),
$M=(S,A,P,R,\gamma)$ \cite{sutton2018,puterman1994}. Cada episódio corresponde à
inspeção de uma imagem. A janela de atenção é representada por:

\begin{equation}
w_t = (c_x, c_y, l, h),
\end{equation}

\noindent em que $c_x$ e $c_y$ são as coordenadas normalizadas do centro da janela, e
$l$ e $h$ sua largura e altura. Os valores são restringidos ao intervalo $[0,1]$, com
$0{,}10 \leq l,h \leq 1{,}00$.

\subsection{Estado e Evidência Visual}

O estado combina informações espaciais da janela, progresso temporal do episódio, valor
atual de interseção sobre união ($IoU_t$), descritores visuais do recorte observado e
probabilidades diagnósticas produzidas por um classificador auxiliar treinado apenas no
conjunto de treino. Esse classificador utiliza estatísticas RGB, intensidade em escala
de cinza, bordas Sobel e histogramas HSV. Sua função é fornecer evidência visual
compacta ao PPO, sem substituir a política de inspeção ativa.

O alinhamento entre a janela atual $B_t$ e as caixas anotadas $G_i$ é medido por:

\begin{equation}
IoU_t = \max_i \frac{|B_t \cap G_i|}{|B_t \cup G_i|}.
\end{equation}

\subsection{Ações e Recompensa}

O espaço de ações é discreto e contém 13 ações: dez ações controlam a janela de
atenção, incluindo deslocamentos, zoom e alteração de largura e altura; as três ações
restantes correspondem às classes diagnósticas. Caso o limite de passos seja atingido
sem uma classificação explícita, a decisão final é definida pela evidência visual da
janela atual.

A recompensa foi proposta para combinar localização, classificação e custo de inspeção:

\begin{equation}
R_t = \alpha \Delta IoU_t + \beta IoU_t - \lambda C_t + R_{cls},
\end{equation}

\noindent em que $\Delta IoU_t = IoU_t - IoU_{t-1}$, $C_t$ representa o custo de
inspeção e $R_{cls}$ corresponde à recompensa terminal de classificação. Predições
corretas recebem recompensa positiva, predições incorretas são penalizadas, e decisões
com maior IoU final recebem bonificação adicional.

Foram avaliadas três variantes: \textit{PPO-localization}, com foco no aumento do IoU;
\textit{PPO-balanced}, que combina localização e classificação; e
\textit{PPO-efficient}, que penaliza janelas grandes e utiliza menor orçamento de
passos.

\subsection{Treinamento e Avaliação}

A política PPO foi implementada como uma MLP com duas camadas ocultas de 64 neurônios.
Os experimentos utilizaram quatro ambientes paralelos, \texttt{n\_steps}=128,
\texttt{batch\_size}=128, \texttt{n\_epochs}=4, \texttt{clip\_range}=0{,}2 e
\texttt{ent\_coef}=0{,}02. Cada variante foi treinada por 50.000 timesteps.

\begin{table}[htbp]
\centering
\caption{Configurações das variantes PPO avaliadas.}
\label{tab:ppo_variants}
\begin{tabular}{lcccc}
\hline
\textbf{Variante} & \textbf{Foco} & \textbf{Passos} & \textbf{Taxa de aprendizado} & $\boldsymbol{\gamma}$ \\
\hline
\textit{PPO-localization} & Localização & 8 & $3{,}0 \times 10^{-4}$ & 0{,}95 \\
\textit{PPO-balanced} & Localização + classificação & 8 & $2{,}5 \times 10^{-4}$ & 0{,}96 \\
\textit{PPO-efficient} & Eficiência de inspeção & 6 & $2{,}5 \times 10^{-4}$ & 0{,}94 \\
\hline
\end{tabular}
\end{table}

A avaliação final foi realizada no conjunto de teste usando acurácia, Macro-F1, taxa de
classificação, melhor IoU médio, IoU final médio, número médio de passos e área média
da janela. A Macro-F1 foi adotada como métrica principal devido ao desbalanceamento das
classes. As variantes foram comparadas com três baselines: classe majoritária,
aleatório uniforme e aleatório estratificado.
